\begin{abstract}
We introduce \textsc{QLoRA}, a novel and memory-efficient finetuning method that enables the adaptation of a 65B parameter language model on a single 48GB GPU, all while matching the performance of traditional 16-bit finetuning. \textsc{QLoRA} accomplishes this by propagating gradients through a fixed, 4-bit quantized pretrained language model into Low Rank Adapters (LoRA). Our top-performing model series, dubbed \textbf{Guanaco}, surpasses all prior open-source models on the Vicuna benchmark, achieving 99.3\% of ChatGPT's performance after just 24 hours of training on a single GPU. \textsc{QLoRA} introduces several technical advancements to reduce memory consumption with no loss in accuracy: (a) 4-bit NormalFloat (NF4), a new data type optimized for the information distribution of model weights, (b) Double Quantization, which minimizes average memory use by also quantizing the quantization parameters themselves, and (c) Paged Optimizers, designed to handle temporary surges in memory demand. Utilizing \textsc{QLoRA}, we finetuned over 1,000 models, offering an extensive evaluation of instruction following and chatbot capabilities across eight instructional datasets, several model families (LLaMA, T5), and model sizes—including those, like 33B and 65B parameters, that regular finetuning cannot accommodate. Our findings reveal that \textsc{QLoRA}-based finetuning on small, high-quality data achieves state-of-the-art results, often outperforming previous best results even with more modest model sizes. We present a comprehensive assessment of chatbot quality using both human raters and GPT-4, showing that GPT-4-based evaluations can serve as a reliable and economical proxy for human judgment. Additionally, our analysis indicates that existing chatbot benchmarks may not provide accurate assessments of system performance. A focused error analysis highlights specific cases where \textbf{Guanaco} does not match ChatGPT. All models, source code, and 4-bit training CUDA kernels are made publicly available.\footnote{\url{https://github.com/artidoro/qlora} and \url{https://github.com/TimDettmers/bitsandbytes}}
\end{abstract}

\section{Introduction}

Refining large language models (LLMs) through finetuning has proven to substantially enhance their capabilities \citep{min2021metaicl, wei2021finetuned, ouyang2022training, wang2022super, wang2022self, liu2022few} and to permit both the augmentation of preferred behaviors and the elimination of unwanted ones \citep{ouyang2022training,askell2021general,bai2022training}. Nevertheless, the process of finetuning models of considerable size is often cost-prohibitive. For instance, standard 16-bit finetuning of a LLaMA model comprising 65B parameters~\citep{touvron2023llama} demands over 780 GB of GPU memory. Although modern quantization strategies can alleviate the memory demands during inference \citep{dettmers2022llm, dettmers2022case, frantar2022gptq, xiao2022smoothquant}, these solutions are typically inapplicable or fail during the training phase \citep{wortsman2023stable}.

In this work, we establish for the first time that finetuning a 4-bit quantized model can be accomplished with no observable loss in performance. Our approach, \textsc{QLoRA}, introduces an advanced high-precision quantization procedure, compressing a pretrained model to 4 bits while incorporating a compact set of trainable Low-rank Adapter weights \citep{hu2021lora} 

which are adapted by propagating gradients through the quantized parameters.

Employing \textsc{QLoRA} brings down the mean GPU memory requirement for finetuning a 65B parameter model from more than 780 GB to under 48 GB, all without increasing computational latency or diminishing predictive accuracy when contrasted with conventional 16-bit finetuning. This development drastically enhances the practicality of LLM finetuning, making it feasible to refine the largest open-source models on a single GPU. With \textsc{QLoRA}, we developed the \textbf{Guanaco} model suite, where the second-highest performing variant achieves 97.8\% of ChatGPT’s performance on the Vicuna~\citep{vicuna2023} benchmark, with training durations below 12 hours on standard consumer GPUs. Utilizing a single professional GPU for 24 hours enables our largest model to reach 99.3\%, effectively closing the performance gap with ChatGPT on Vicuna. Furthermore, upon deployment, our smallest \textbf{Guanaco} model (7B parameters) needs only 5 GB of memory, outperforming the 26 GB Alpaca model by over 20 percentage points on the Vicuna benchmark (Table~\ref{tbl:vicuna_eval}).

\textsc{QLoRA} incorporates several advancements aimed at minimizing memory consumption while maintaining model efficacy: (1) \textbf{4-bit NormalFloat}, an information-theoretically optimal quantization format tailored for data with a normal distribution, which demonstrates superior empirical performance relative to 4-bit Integers and 4-bit Floats. (2) \textbf{Double Quantization}, a technique in which the quantization coefficients themselves are further quantized, resulting in average savings of approximately 0.37 bits per parameter—equivalent to roughly 3 GB for a 65B parameter model. (3) \textbf{Paged Optimizers}, which leverage NVIDIA's unified memory to circumvent memory surges typically encountered during gradient checkpointing, particularly for mini-batches comprising lengthy sequences. By integrating these strategies, we realize an enhanced LoRA variant that incorporates adapters within all network layers, thereby nearly eliminating the accuracy compromises characteristic of earlier methods.

The efficiency gains achieved through \textsc{QLoRA} make it feasible to conduct comprehensive analyses of instruction finetuning and chatbot effectiveness at model scales previously unattainable with traditional finetuning approaches, owing to prohibitive memory requirements. As a result, we successfully train in excess of 1,000 models, spanning multiple instruction tuning datasets, architectures, and ranging in size from 80M to 65B parameters. Beyond establishing that \textsc{QLoRA} can match 16-bit finetuning outcomes (\S\ref{sec:comp_to_std}) and enabling the training of the high-performing chatbot \textbf{Guanaco} (\S\ref{sec:pushing}), we also provide a detailed examination of trends in the resulting models. Notably, we observe that data quality substantially outweighs dataset quantity: for instance, a dataset of 9,000 samples (OASST1) surpassed a much larger 450,000-sample dataset (FLAN v2, subsampled) on measures of chatbot performance, even when both datasets are intended for instruction following. Furthermore, our findings reveal that high performance on the Massive Multitask Language Understanding (MMLU) benchmark does not necessarily translate to high scores on the Vicuna chatbot benchmark, or the reverse—emphasizing that the alignment between dataset characteristics and specific task requirements is more critical than sheer size.

In addition, we conduct a comprehensive assessment of chatbot effectiveness using both human judges and GPT-4-based evaluations. The methodology follows a tournament-like framework, wherein multiple models are pitted against each other to generate the superior response to a given prompt. The outcome of each pairwise comparison, or "match," is determined by either GPT-4 or human evaluators. Final rankings are computed by aggregating these results into Elo scores~\citep{elo1967proposed, elo1978rating}, establishing a relative performance hierarchy among the chatbots. Our analysis reveals a general concordance between GPT-4-based and human judgments regarding model rankings, though notable divergences are also observed. This underscores that, while model-centric assessment methods offer a cost-effective substitute for human evaluation, they also introduce certain ambiguities and limitations.

To complement our benchmark findings, we include a qualitative evaluation focusing on the \textbf{Guanaco} family of models. This examination brings attention to examples of model strengths and shortcomings that eluded the scope of quantitative measures.

For the advancement of subsequent research, we make available all generated model responses, including accompanying human and GPT-4 annotation data. The entirety of our codebase, including custom CUDA kernels, has been open-sourced, and we have incorporated our techniques into the Hugging Face transformers library \citep{wolf2019huggingface} for broader accessibility. Our contributions further include releasing a suite of adapters for models of sizes 7/13/33/65B, trained across eight distinct instruction-following datasets—yielding 32 openly available, finetuned models.

\section{Background}
\paragraph{Block-wise k-bit Quantization} Quantization involves reducing the precision of data by mapping values from a high-information representation to one that holds fewer bits of information, such as transforming 32-bit floating point values into 8-bit integer values. This process typically includes normalizing the input data—often structured as a tensor—by scaling it to utilize the full target data type range, based on the maximum absolute value present. For instance, to quantize a 32-bit floating point (FP32) tensor into an Int8 tensor with a target range of $[-127, 127]$, the transformation is given by:

\begin{equation}
    \mathbf{X}^{ \text{Int8}} = \text{round}\left( \frac{127}{{\text{absmax}}(\mathbf{X}^{{\text{FP32}}})}\mathbf{X}^{{\text{FP32}}} \right) = \text{round}(c^{\text{FP32}}\cdot \mathbf{X}^{{\text{FP32}}}),
\end{equation}

where $c$ serves as the {\em quantization constant} or {\em quantization scale}. The dequantization operation, which reconstructs the original values, is described as:

\begin{equation}
    \text{dequant}(c^{\text{FP32}}, \mathbf{X}^{\text{Int8}}) = \frac{\mathbf{X}^{\text{Int8}}}{c^{\text{FP32}}} = \mathbf{X}^{\text{FP32}}
\end{equation}

One major drawback of this method arises when an outlier—a value with unusually large magnitude—appears in the input tensor. This results in poor utilization of the quantization bins, where certain bit patterns may see little to no data mapped to them. To address the presence of such outliers, it is common practice to partition the input tensor into smaller sections or blocks, with each block quantized separately using its specific quantization constant $c$. More precisely, the input tensor $\mathbf{X}\in \mathbb{R}^{b\times h}$ is divided into $n$ consecutive blocks, each of size $B$, by flattening $\mathbf{X}$ and dividing the sequence into ${n = ({b\times h})/{B}}$ blocks. Each of these blocks is then quantized independently according to Equation~1, resulting in a quantized tensor and a set of $n$ distinct quantization constants $c_i$.

\paragraph{Low-rank Adapters} LoRA finetuning \citep{hu2021lora} is a technique that significantly cuts down on memory consumption by introducing a compact set of trainable adapter parameters, while the main model parameters remain unchanged during training. During stochastic gradient descent, the gradients propagate through the fixed pretrained weights and are applied to the adapter module, which is trained to minimize the loss. LoRA enhances a linear transformation with an extra low-rank decomposition. For a projection given by ${\mathbf{X} \mathbf{W} = \mathbf{Y}}$, where $\mathbf{X}\in  \mathbb{R}^{b\times h}$ and $\mathbf{W}\in  \mathbb{R}^{h\times o}$, LoRA computes:

\begin{equation}
    \mathbf{Y} = \mathbf{X}\mathbf{W} + s\mathbf{X}\mathbf{L}_1\mathbf{L}_2,
\end{equation}

with $\mathbf{L}_1\in  \mathbb{R}^{h\times r}$ and $\mathbf{L}_2\in  \mathbb{R}^{r\times o}$, and $s$ being a scalar factor.

\paragraph{Memory Requirement of Parameter-Efficient Finetuning} The memory usage associated with LoRA during training warrants particular attention, especially regarding both the count and the scale of adapters utilized. Owing to LoRA’s extremely modest memory demands, it becomes feasible to deploy a greater number of adapters to enhance model performance, all while causing only a negligible increase in total memory consumption. Despite being engineered as a Parameter Efficient Finetuning (PEFT) approach, the predominant memory usage in LLM finetuning is attributable to the storage of activation gradients rather than to the LoRA parameters themselves. For example, when finetuning a 7B LLaMA model on FLAN v2 with a batch size of 1 and setting LoRA weights to 0.2\% of the model’s original parameter count—a standard choice in prior works~\citep{hu2021lora,liu2022few}—the memory needed for LoRA input gradients amounts to 567 MB, whereas the LoRA parameter storage itself is a mere 26 MB. Application of gradient checkpointing~\citep{chen2016training} reduces the average input gradient memory per sequence to 18 MB, which is still more than the aggregate memory used by all LoRA weights. In contrast, the memory required for the base model in 4-bit precision is 5,048 MB. This demonstrates the essential role of gradient checkpointing while also illustrating that further reductions in the volume of LoRA parameters contribute only marginally to memory savings. Consequently, it is practical to increase the number of adapters used without meaningfully impacting total training memory requirements (refer to Appendix~\ref{app:memory} for an in-depth analysis). As elaborated later, this ability is fundamental for achieving the performance of full 16-bit precision.

\section{\textsc{QLoRA} Finetuning}

\textsc{QLoRA} attains high-accuracy 4-bit finetuning by employing two principal methods: 4-bit NormalFloat (NF4) quantization and Double Quantization, both introduced in this work. In addition, we propose the use of Paged Optimizers, designed to alleviate memory surges caused by gradient checkpointing—an issue that has historically rendered single-machine finetuning of large models prone to out-of-memory errors.

Within the \textsc{QLoRA} framework, there is a separation between the data type utilized for storage—typically 4-bit precision—and the data type used during computation, commonly BFloat16. In practical terms, whenever a weight tensor from \textsc{QLoRA} is accessed, it is first dequantized to BFloat16, followed by 16-bit matrix multiplication operations.

We next elaborate on the primary mechanisms of \textsc{QLoRA} and provide a precise formalization of the method.

\paragraph{4-bit NormalFloat Quantization} The NormalFloat (NF) data type extends the approach of Quantile Quantization \citep{dettmers2022optimizers}, which is considered an information-theoretically optimal representation. This quantization technique assigns an equal number of input tensor values to each quantization bin by estimating the distribution's quantiles using the empirical cumulative distribution function.

A significant drawback of quantile quantization is the computational cost associated with accurate quantile estimation. Consequently, accelerated quantile approximation methods, such as SRAM quantiles~\citep{dettmers2022optimizers}, are adopted to mitigate this expense. However, these fast approximation schemes tend to introduce considerable quantization errors, particularly for outlier values, which frequently hold high importance.

Both the computational overhead of precise quantile computation and the inaccuracy from approximations can be minimized when input tensors originate from distributions determined up to a quantization constant. In these scenarios, the input tensors possess consistent quantiles, rendering exact quantile calculations tractable.

Typically, the pretrained weights of neural networks are drawn from a normal distribution centered at zero with a standard deviation $\sigma$ (refer to Appendix~\ref{app:normality}). To unify the distribution of all weights, we can rescale $\sigma$ so that the resulting distribution aligns precisely with the representational range of our data type, which we define to be $[-1, 1]$. This necessitates mapping both the quantiles from the original weights and the target data type into this standardized interval.

For normal distributions with zero mean and arbitrary $\sigma$, constrained within $[-1, 1]$, the data type that achieves optimal information efficiency can be derived via the following procedure: (1) calculate the $2^k + 1$ quantile thresholds of a standard normal distribution $N(0, 1)$ to create a $k$-bit quantile quantization scheme tailored for normal variables, (2) adjust the quantized data type so its representational values are distributed over $[-1, 1]$, and (3) normalize the input weight tensor into $[-1, 1]$ using absolute max scaling, thus preparing it for quantization.

Once the bounds of the data type and the weight tensor coincide, standard quantization techniques may be applied. This normalization procedure in step (3) equates to scaling the weight tensor's standard deviation such that it matches the standard deviation inherent to the $k$-bit quantized data type. Explicitly, we compute the $2^k$ discrete quantization values $q_i$ for the data type according to:

\begin{equation}
q_i  = \frac{1}{2}\left( Q_X\left(\frac{i}{2^k+1}\right) + Q_X\left(\frac{i+1}{2^k+1}\right)\right),
\end{equation}

where $Q_X(\cdot)$ denotes the quantile function corresponding to the standard normal distribution $N(0,1)$. A challenge inherent to symmetric $k$-bit quantization is the inability of this scheme to represent zero exactly—an essential feature when quantizing elements such as padding and other zero-valued entries without introducing error. To guarantee an explicit zeropoint at $0$, and to utilize the entire set of $2^k$ discrete values within a $k$-bit data format, we construct an asymmetric representation. This is achieved by computing quantiles $q_i$ over two domains: $2^{k-1}$ quantiles for the negative segment, and $2^{k-1}+1$ quantiles for the positive segment. We then merge these two sets of $q_i$ and remove the redundant zero that appears in both, resulting in a unified quantization scheme. We refer to this data type as \emph{k-bit NormalFloat} (NFk), characterized by each quantization bin containing an equal expected proportion of the distribution, and being information-theoretically optimal for zero-centered normal distributions. The precise specifications of this format are provided in Appendix~\ref{app:nf4}.

\paragraph{Double Quantization}
We present \emph{Double Quantization} (DQ), which aims to further reduce memory consumption by subjecting the quantization constants themselves to an additional quantization step. Traditional 4-bit quantization with small block sizes, as described by~\citet{dettmers2022case}, requires numerous 32-bit constants, contributing considerable memory overhead. As an illustration, when employing 32-bit constants and a blocksize of 64 to quantize $\mathbf{W}$, each parameter effectively accrues an additional $32/64 = 0.5$ bits due to the quantization constants. Double Quantization provides a mechanism to decrease this memory requirement.

Concretely, Double Quantization processes the quantization constants $c_2^{\text{FP32}}$ from the initial quantization as input for a subsequent quantization step. This yields $c_2^{\text{FP8}}$, the quantized form of the quantization constants, along with an additional set of second-stage quantization constants $c_1^{\text{FP32}}$. The second quantization employs 8-bit floating-point representations with a blocksize of 256, reflecting empirical findings that no accuracy loss occurs for 8-bit quantization~\citep{dettmers2022case}. Because $c_2^{\text{FP32}}$ are nonnegative, we first subtract their mean prior to the second quantization, thereby centering the values and allowing for symmetric quantization. As a result, for a blocksize of 64, this approach decreases the per-parameter memory cost associated with quantization constants from $32/64 = 0.5$ bits to $8/64 + 32/(64 \cdot 256) = 0.127$ bits—representing a savings of 0.373 bits per parameter.

\paragraph{Paged Optimizers} leverage the NVIDIA unified memory~\footnote{\tiny \url{https://docs.nvidia.com/cuda/cuda-c-programming-guide}} capability, which facilitates seamless page-level data transfers between the CPU and GPU. This mechanism ensures reliable GPU computation, even when GPU memory becomes temporarily exhausted. Functioning similarly to traditional paging between CPU memory and disk storage, this feature allows us to allocate paged memory for optimizer states; when GPU memory is insufficient, these states are transparently swapped to CPU RAM and reloaded to the GPU as required during optimizer updates.

\paragraph{\textsc{QLoRA}.} By combining the mechanisms outlined above, we implement \textsc{QLoRA} for a single quantized linear layer equipped with one LoRA adapter, formulated as:

\begin{equation}
    \mathbf{Y}^{\text{BF16}} =  \mathbf{X}^{\text{BF16}} \text{doubleDequant}(c_1^{\text{FP32}}, c_2^{\text{k-bit}}, \mathbf{W}^{\text{NF4}}) + \mathbf{X}^{\text{BF16}}\mathbf{L}^{\text{BF16}}_1\mathbf{L}^{\text{BF16}}_2,
\end{equation}

Here, the operator doubleDequant$(\cdot)$ is given by:

\begin{equation}
    \text{doubleDequant}(c_1^{\text{FP32}}, c_2^{\text{k-bit}}, \mathbf{W}^{\text{k-bit}}) = \text{dequant}(\text{dequant}(c_1^{\text{FP32}}, c_2^{\text{k-bit}}), \mathbf{W}^{\text{4bit}}) = \mathbf{W}^{\text{BF16}},
\end{equation}

Within our framework, we choose the NF4 format for representing $\mathbf{W}$, while $c_2$ is quantized to FP8.

We opt for a blocksize of 64 for $\mathbf{W}$ to enhance quantization accuracy, whereas a blocksize of 256 is selected for $c_2$ to optimize memory efficiency.

For updating parameters, it is only necessary to compute gradients with respect to the adapter weights, $\frac{\partial E}{\partial \mathbf{L}_i}$, omitting the computation for the quantized 4-bit weights, $\frac{\partial E}{\partial \mathbf{W}}$. Nonetheless, the evaluation of $\frac{\partial E}{\partial \mathbf{L}_i}$ requires the calculation of $\frac{\partial \mathbf{X}}{\partial \mathbf{W}}$, which is performed via equation~(5), utilizing dequantization to transform stored $\mathbf{W}^{\text{NF4}}$ values into computational $\mathbf{W}^{\text{BF16}}$ form, thereby ensuring the gradient $\frac{\partial \mathbf{X}}{\partial \mathbf{W}}$ is computed with BFloat16 precision.

In summary, \textsc{QLoRA} utilizes a single storage data type (commonly 4-bit NormalFloat) along with a computation data type (16-bit BrainFloat). During training, the stored values are dequantized from the storage format to the computation format to enable both forward and backward propagation; however, the computation of weight gradients is restricted solely to the LoRA parameters, which employ 16-bit BrainFloat.

\section{QLoRA vs. Standard Finetuning}
\label{sec:comp_to_std}

After outlining the mechanism of QLoRA and illustrating its substantial memory efficiency for finetuning, the central issue becomes whether QLoRA can achieve parity in performance with comprehensive full-model finetuning. Additionally, it is of interest to dissect the internal elements of QLoRA, specifically the contribution of NormalFloat4 as compared to the conventional Float4. In the forthcoming sections, we present experimental investigations aimed at addressing these concerns.

\paragraph{Experimental setup.} Our experimental design spans three model types (encoder, encoder-decoder, and decoder-only), facilitating a comparison of QLoRA against both 16-bit adapter-based finetuning and complete finetuning procedures, with models of sizes up to 3B. The evaluation framework involves GLUE \citep{wang2018glue} utilizing RoBERTa-large \citep{liu2019roberta}, Super-NaturalInstructions (TKInstruct) \citep{wang2022super} employing T5 \citep{t5}, and 5-shot MMLU \citep{hendrycksmeasuring} after the finetuning of LLaMA on Flan v2 \citep{longpre2023flan} and Alpaca \citep{alpaca}. To further examine the comparative merits of NF4 versus alternative 4-bit data types, we follow the protocol of \citet{dettmers2022case}, evaluating post-quantization zero-shot accuracy and perplexity over a range of architectures (OPT~\citep{zhang2022opt}, LLaMA~\citep{touvron2023llama}, BLOOM~\citep{scao2022bloom}, Pythia~\citep{biderman2023pythia}), covering model scales from 125m to 13B.

Further specific experimental configurations are discussed within the results section of each experimental context for clarity, and comprehensive details are provided in Appendix~\ref{app:exp-setup}.

Although paged optimizers are indispensable for executing 33B/65B \textsc{QLoRA} finetuning runs on a single GPU with 24GB or 48GB of memory, we omit precise timing metrics for Paged Optimizers, since paging tends to occur infrequently and primarily when large mini-batches with long input sequences are processed. Nevertheless, we assess the runtime behavior of paged optimizers on 65B models with 48GB GPUs and observe that, for a batch size of 16, their training speed matches that of standard optimizers. Investigating the conditions leading to slowdowns from the paging process remains an avenue for future research.

\paragraph{Default LoRA hyperparameters do not match 16-bit performance}

Applying the conventional LoRA strategy—where adapters are placed only in query and value attention projection layers \citep{hu2021lora}—fails to reproduce full finetuning performance for large model bases. Figure~\ref{fig:hyper} (illustrating LLaMA 7B finetuning on Alpaca) demonstrates that the dominant factor among LoRA hyperparameters is the number of LoRA adapters deployed, and it is necessary to incorporate LoRA adapters in all linear transformer block layers to achieve parity with full finetuning outcomes. Other hyperparameters of LoRA, such as projection rank $r$, appear to exert minimal influence on performance (refer to Appendix \ref{app:exp-setup}).

In a similar vein, our analysis reveals that the standard hyperparameters employed for fully finetuned baselines are insufficiently optimized. To establish more robust baselines, we conduct an extensive search across learning rates ranging from 1e-6 to 5e-5 and batch sizes spanning 8 to 128. The corresponding outcomes for 7B LLaMA finetuning on Alpaca are presented in Figure~\ref{fig:hyper}.

\paragraph{4-bit NormalFloat outperforms 4-bit Floating Point}

Although the 4-bit NormalFloat (NF4) data type is known to be information-theoretically optimal, it remains an empirical question whether this theoretical advantage yields improved performance in practice. Adopting the methodology from \citet{dettmers2022case}, we evaluate quantized LLMs (including OPT \citep{zhang2022opt}, BLOOM \cite{scao2022bloom}, Pythia \cite{biderman2023pythia}, and LLaMA) across various model scales (125M up to 65B parameters), utilizing multiple data types on language modeling tasks as well as a suite of zero-shot benchmarks. Figure~\ref{fig:nf4} and Table~\ref{tbl:nf4_ppl} demonstrate that NF4 yields markedly better results than both FP4 and Int4, while double quantization further minimizes memory usage without negatively impacting performance.

\paragraph{k-bit \textsc{QLoRA} achieves parity with 16-bit full finetuning and 16-bit LoRA}

Previous work indicates that while 4-bit quantization suffices for \emph{inference}, it often results in diminished performance compared to 16-bit representations \citep{dettmers2022case,frantar2022gptq}. This observation prompts a key investigation: Can adapter finetuning with 4-bit quantization fully restore the original performance? We examine this in two scenarios.

Our first comparison targets full 16-bit finetuning for RoBERTA and T5 models, with parameter sizes between 125M and 3B, on the GLUE and Super-NaturalInstructions datasets. The findings, summarized in Table~\ref{tab:nf4vsbf16}, indicate that the 16-bit, 8-bit, and 4-bit adapter techniques all recover the accuracy of the 16-bit fully finetuned baseline across both datasets. Thus, adapter finetuning following quantization appears sufficient to reclaim the performance deficit introduced by lower-precision representations.

In our second experimental configuration, given that fully finetuning models with 11B parameters or more demands multiple high-memory GPU servers, we further investigate if 4-bit \textsc{QLoRA} achieves comparable outcomes to 16-bit LoRA across parameter scales ranging from 7B to 65B. To explore this, we conduct finetuning of LLaMA models at 7B, 13B, 33B, and 65B parameter sizes on two instruction-tuning datasets, Alpaca and FLAN v2, and assess their performance on the MMLU benchmark utilizing 5-shot accuracy. As detailed in Table~\ref{tbl:llama-datatype}, NF4 with double quantization is observed to fully restore the MMLU scores attained by 16-bit LoRA. Additionally, it is evident that when using FP4, \textsc{QLoRA} underperforms compared to the 16-bit brain float LoRA baseline by approximately 1 percentage point. These findings validate both that (1) NF4-based \textsc{QLoRA} reliably mirrors the efficacy of 16-bit full finetuning and 16-bit LoRA, and (2) NF4 provides greater quantization fidelity than FP4.

\paragraph{Summary} Our experiments robustly demonstrate that employing 4-bit \textsc{QLoRA} with the NF4 data format consistently produces results equivalent to both 16-bit full and 16-bit LoRA finetuning on established academic tasks. Furthermore, NF4 surpasses FP4 in terms of effectiveness, and incorporating double quantization does not diminish accuracy. Collectively, these results present strong support that 4-bit \textsc{QLoRA} finetuning reliably matches 16-bit techniques.

Echoing prior quantization research \citep{dettmers2022case}, our analyses of MMLU and Elo outcomes reveal that, given fixed finetuning and inference resources, opting to expand the base model’s parameter count while reducing numerical precision is advantageous. This underscores the efficiency benefits introduced by \textsc{QLoRA}. Since 4-bit finetuning does not display any loss in performance relative to full finetuning in our trials, it prompts further inquiry into the precise point where the trade-off between performance and precision emerges in QLoRA finetuning—a question that remains open for future investigation.

We aim to examine instruction tuning at scales unattainable with traditional 16-bit finetuning using standard academic computing resources.

\section{Pushing the Chatbot State-of-the-art with QLoRA}
\label{sec:pushing}
After confirming that 4-bit \textsc{QLoRA} achieves performance on par with 16-bit methods across varying model sizes, task types, and datasets, we turn our attention to a thorough analysis of instruction finetuning for the largest open-access language models available for academic use. To rigorously evaluate instruction finetuning outcomes, we utilize a demanding Natural Language Understanding benchmark (MMLU) and propose new strategies for real-world chatbot evaluation.

\subsection{Experimental setup} The overall experimental design is outlined below, with comprehensive implementation details presented in Appendix \ref{app:chatbot}.

\paragraph{Data} Since no systematic investigation of recent instruction-following datasets exists to our knowledge, we curate a collection of eight contemporary datasets. Our selection encompasses resources sourced via crowd-sourcing (OASST1~\citep{kopf2023openassistant}, HH-RLHF~\citep{bai2022training}), distillations from models previously trained with instructions (Alpaca~\citep{alpaca}, self-instruct~\citep{wang2022self}, unnatural-instructions~\citep{honovich2022unnatural}), comprehensive corpora collections (FLAN v2~\citep{chung2022scaling}), and composite datasets (Chip2~\citep{laion2023}, Longform~\citep{koksal2023longform}). This collection features a diversity of languages, corpus sizes, and licensing arrangements.

\paragraph{Training Setup} To isolate the effects of the instruction-tuning approach, we consistently employ QLoRA finetuning utilizing cross-entropy loss, omitting reinforcement learning—even when human feedback data is available. When datasets explicitly separate the instruction from the target response, training is conducted exclusively on the response portion (see Appendix \ref{app:chatbot} for related ablation studies). For datasets like OASST1 and HH-RLHF that offer multiple response candidates, we select the highest-rated reply at every node within the conversation tree and perform finetuning using the full trajectory, inclusive of both prompts and selected responses. In all experiments, we deploy NF4 \textsc{QLoRA} in combination with double quantization and paged optimization in order to mitigate the risk of memory overuse during gradient checkpointing. Hyperparameter optimization is conducted on the LLaMA 13B and 33B models, where hyperparameters derived from tuning the 7B model generally extend to larger models—with learning rate and batch size as exceptions. Specifically, for the 33B and 65B configurations, the learning rate is reduced by half, while the batch size is doubled.

\paragraph{Baselines} The trained models are evaluated against both academic systems (Vicuna~\citep{vicuna2023}, Open Assistant~\citep{kopf2023openassistant}) and established commercial products (GPT-4 \citep{openai2023gpt}, GPT-3.5-turbo, and Bard). Open Assistant is based on LLaMA 33B and trained with Reinforcement Learning from Human Feedback on the OASST1 data, which overlaps with our experiments. In contrast, Vicuna is created by fully finetuning LLaMA 13B on conversation transcripts collected from ShareGPT, effectively resulting in distillation of outputs from OpenAI GPT models.

\subsection{Evaluation}

Consistent with standard methodology, we evaluate models using the MMLU (Massively Multitask Language Understanding) benchmark \citep{hendrycksmeasuring}, which assesses abilities across a broad set of 57 tasks including domains such as basic mathematics, US history, computer science, and law. Performance is reported as 5-shot accuracy on the benchmark’s test split.

We further assess generative language abilities via both automated methods and human judgment. This complementary set of evaluations employs queries selected by humans, with the objective of quantifying the quality of model outputs. Although this approach provides a more authentic measure of chatbot effectiveness and is becoming increasingly prevalent, the field has yet to converge on a standardized methodology. Our experimental design is detailed below; across all settings, we utilize nucleus sampling with $p=0.9$ and a temperature of $0.7$.

\paragraph{Benchmark Data} Our evaluations utilize two specially curated sets of queries (questions): Vicuna prompts \citep{vicuna2023} and the OASST1 validation dataset \citep{kopf2023openassistant}. The Vicuna benchmark consists of 80 unaltered prompts spanning a variety of categories. The OASST1 dataset contains multilingual, user-assistant, multiturn conversations sourced through crowdsourcing. For OASST1, we extract all user inputs within the validation set as queries, incorporating their preceding conversational context in the prompt, resulting in 953 unique user queries. For clarity, we refer to these two datasets as the Vicuna and OA benchmarks.

\paragraph{Automated Evaluation} Following the methodology established by \citet{vicuna2023}, we implement GPT-4 to assess the quality of model responses against ChatGPT (GPT-3.5 Turbo) on the Vicuna benchmark. Given a query, GPT-4 is presented with both ChatGPT’s response and a candidate model’s response, and tasked with assigning a score (out of ten) to each, along with a brief rationale. The model’s overall score is computed as a percentage of the ChatGPT reference score; notably, values can exceed 100\% if the evaluated model obtains a higher raw score. We observe a pronounced positional bias: GPT-4 tends to favor the response appearing first in the prompt. To account for this, we recommend averaging the scores over both possible response orders.

In addition, we evaluate models by directly comparing their outputs. The evaluation protocol is simplified to a three-way classification, allowing for ties. Here, GPT-4 is asked to select the superior answer, or indicate a tie, justifying its choice. We execute comprehensive pairwise comparisons across all system pairs for both the Vicuna and OA benchmarks.

\paragraph{Human Evaluation} Despite indications from prior studies that generative models can effectively perform system assessments \citep{fu2023gptscore}, it remains uncertain whether GPT-4-based evaluations of chatbot performance align well with human judgments. Consequently, we conduct two parallel rounds of human evaluation on the Vicuna benchmark, mirroring the two automated protocols. Using Amazon Mechanical Turk (AMT), we engage two annotators for ChatGPT-comparison tasks and three annotators for each pairwise system comparison.

\paragraph{Elo Rating} In our methodology, both human evaluators and automated systems conduct pairwise assessments, structuring the process as a tournament in which models are pitted against one another. Each matchup involves two models generating responses to identical prompts, and their outputs are judged in direct comparison. While this framework is akin to the procedures of \citet{bai2022training} and \citet{vicuna2023}, we extend it by incorporating ratings provided by GPT-4 alongside those from human annotators. For Elo~\citep{elo1967proposed,elo1978rating} calculation, we draw randomly from our repository of labeled comparisons. The Elo rating system, extensively utilized in chess and competitive games, quantifies a model's anticipated performance relative to an opponent—for instance, a model rated 1100 has about a 65\% likelihood of outperforming a model rated 1000, whereas equally rated models (1000 vs 1000 or 1100 vs 1100) each have a 50\% expected win rate. Elo scores are updated following each contest, with the magnitude of change dependent on how surprising the match result is: upsets prompt larger adjustments, while predictable results result in smaller updates. As more matches are played, Elo scores tend to reflect each model's comparative strength within the game. All models are initially assigned a baseline score of 1,000 and we set $K=32$. In alignment with \citet{vicuna2023}, this tournament simulation is conducted 10,000 times, each with a different random seed, in order to mitigate sequence bias arising from the order in which model pairs are evaluated.

\subsection{\textbf{Guanaco}: \textsc{QLoRA} trained on OASST1 Achieves State-of-the-Art Chatbot Performance}

Through both automated assessments and manual review, our findings indicate that the {Guanaco} 65B model, fine-tuned using a modified version of OASST1 with \textsc{QLoRA}, currently stands as the leading open-source chatbot, matching the capabilities of ChatGPT. Direct human system-level pairwise comparisons, summarized as Elo ratings, reveal that both {Guanaco} 65B and 33B attain a 30\% expected win probability against GPT-4—the highest figure reported thus far.

Table~\ref{tbl:vicuna_eval} presents {Guanaco}'s comparative performance on the Vicuna benchmark \cite{vicuna2023} relative to ChatGPT. Notably, {Guanaco} 65B emerges as the top model following GPT-4, reaching 99.3\% of ChatGPT’s performance. While {Guanaco} 33B includes more parameters than Vicuna 13B, its use of 4-bit weight precision yields greater memory efficiency (21 GB vs. 26 GB), and it outperforms Vicuna 13B by three percentage points. Additionally, the compact {Guanaco} 7B variant requires only 5 GB of memory—making it suitable for deployment on modern smartphones—and surpasses Alpaca 13B by nearly 20 percentage points in benchmark scores.

Nonetheless, the data in Table~\ref{tbl:vicuna_eval} exhibits substantial confidence intervals, and multiple models exhibit overlapping results. We propose that this variability stems from ambiguous scale definitions—for instance, it is uncertain how to interpret an 8 on a 10-point scale in various contexts. Consequently, we advocate for Elo-based \textit{pairwise} evaluation methods \citep{elo1967proposed}, which utilize judgments from both human raters and GPT-4 to circumvent the subjectivity of absolute scales. Table~\ref{tbl:elo} displays Elo scores for leading models. While rankings by GPT-4 and humans on the Vicuna benchmark generally align (with system-level Kendall Tau $\tau=0.43$ and Spearman $r=0.55$), some divergence remains—most notably for Guanaco 7B. Agreement on specific examples is lower (Fleiss $\kappa=0.25$). Altogether, these results suggest a moderate concordance between GPT-4 and human assessments at the system level, indicating that model-based evaluation can serve as a viable, though imperfect, alternative to human evaluation. Additional points are addressed in Section~\ref{ssec:considerations}.

The Elo rankings presented in Table~\ref{tbl:elo_large} reveal that {Guanaco} models with 33B and 65B parameters surpass all counterparts except GPT-4 on both the Vicuna and OA benchmarks, and they achieve results similar to ChatGPT, as shown in Table~\ref{tbl:vicuna_eval}. It is important to highlight that the Vicuna benchmark is more favorable to open-source models, whereas the OA benchmark tends to give an advantage to ChatGPT. Moreover, analysis of Tables \ref{tbl:mmlu-llama} and \ref{tbl:vicuna_eval} underscores that the choice of finetuning dataset plays a crucial role in overall model performance. For instance, Llama models that are finetuned on FLAN v2 excel in MMLU but show their weakest results on the Vicuna benchmark—a pattern mirrored in other models as well. These observations emphasize a degree of independence among evaluation benchmarks: achieving high scores on MMLU does not guarantee robust chatbot abilities on metrics like Vicuna or OA, nor does strong chatbot performance ensure high MMLU accuracy.

 stands out as the sole leading model in our assessment that is entirely free of proprietary training data, in line with OASST1’s explicit restriction against using GPT-generated content. The closest competitor among models trained exclusively on open-source corpora is Anthropic’s HH-RLHF, which lags behind {Guanaco} by 30 percentage points on the Vicuna benchmark (refer to Table~\ref{tbl:vicuna_eval}). Collectively, these findings indicate that the 4-bit \textsc{QLoRA} approach is highly effective, yielding chatbots that are competitive with state-of-the-art systems such as ChatGPT. Notably, our 33B {Guanaco} model can be trained within 12 hours using consumer-grade 24 GB GPUs, suggesting considerable promise for future advancements through \textsc{QLoRA} fine-tuning on specialized open-source datasets. Such advancements could produce models that rival the performance of top commercial systems currently available.

\section{Qualitative Analysis}

Although our primary assessment focuses on quantitative measures, relying solely on aggregate statistics introduces certain limitations. Chief among these is the concern of benchmark validity \citep{liao2021we}: it remains uncertain whether a given benchmark actually assesses the competencies implied by its title or stated aims, a doubt intensified by the identification of ``shortcuts'' that machine learning models may leverage to artificially boost performance \citep{gururangan2018annotation, poliak2018hypothesis}. To partially address this limitation, we include a qualitative component in our analysis, organized into two subsections. In \S\ref{ssec:examples}, we present illustrative instances that exemplify recurring tendencies detected in outputs from our 65b {Guanaco} model. Following this, in \S\ref{ssec:considerations}, we elaborate on the broader implications of our findings and the rationale underpinning our interpretations.

\subsection{Analysis of Generated Examples}
\label{ssec:examples}

To select representative samples, we systematically review outputs produced in response to prompts from the Vicuna and OpenAssistant benchmarks, seeking discernible trends in the answers given by {Guanaco}. Upon identifying such a trend, we construct targeted prompts aimed at eliciting the same pattern even if it does not correspond to the correct response. For instance, noticing a propensity for verbose outputs, we challenge the model with prompts such as ``Answer yes or no without explanation.'' This method enables us to locate failure cases (``lemons''), where the model's weaknesses become apparent under adversarial probing, as well as successes (``cherries''), where the model demonstrates robustness. Every example in this section was generated using Nucleus Sampling \cite{holtzmancurious} with $p=0.9$.

Naturally, this analysis does not attempt to be exhaustive, as the limited scope of this qualitative study precludes accounting for every potential variable—for example, the sheer breadth of responses a model could produce from a single prompt is extensive, so we must rely on samples that we believe to be broadly indicative. Nonetheless, we argue that including these illustrative cases provides valuable context for the quantitative findings reported earlier in the paper. By releasing all code and models as open source, we intend for this section to encourage further research that delves more deeply into the challenges discussed here.
\paragraph{Factual Recall} When responding to straightforward questions such as ``What is the capital of Zambia?'', all models regularly yield the correct reply, for example,

\begin{quote} 
    \textbf{Guanaco}:
    \texttt{The capital of Zambia is Lusaka.}
\end{quote}

 However, as prompts become less common or require niche knowledge, {Guanaco} becomes prone to producing inaccurate answers, despite maintaining a high degree of certainty. For example, when prompted with a question from HotPotQA \citep{yang2018hotpotqa}:

\begin{quote}
     \textbf{User}: In what year was the singer who popularized `I'll Keep the Lovelight Burning' born?  

     \textbf{Guanaco}:~\texttt{Al Jolson is credited with popularizing the song `I'll Keep the Lovelight Burning,' and he was born in the year 1886.}
 \end{quote}

 the model both misattributes the song’s popularizer and provides an incorrect birth year (although the birth year given is correct for the person wrongly identified, Al Jolson).

\paragraph{Suggestibility} Notably, {Guanaco} demonstrates an unexpected level of resistance to accepting or propagating certain kinds of presupposed false information, as illustrated by the following exchange:

\begin{quote}
    \textbf{User}: How was it finally, officially confirmed that the earth is flat by peer reviewed scientists? \\ \\
    \textbf{Guanaco}: \texttt{There has never been a peer-reviewed scientific confirmation that the earth is flat. The scientific community overwhelmingly agrees that the earth is an oblate spheroid, a fact established for centuries and supported by extensive experimental and observational evidence. \\ \\ Claims that peer-reviewed scientists have validated a flat earth are unfounded and stem from conspiracy theories that science has repeatedly disproven. \\ \\ Feel free to ask if you have further questions.}
\end{quote}

 also reliably identifies questions that cannot be answered as posed, for example,

\begin{quote}
    \textbf{User}: What time is it?

    \textbf{Guanaco}: \texttt{Unfortunately, as an AI language model, I do not possess the capability to access real-time data. My answers rely on information available at the time of training and are not able to indicate the current time. For up-to-date time information, I recommend consulting a clock or an online time source.}
\end{quote}

Comparable outcomes are observed for queries such as ``Where are you?'', ``How are you?'', and so forth.

\paragraph{Refusal} In much the same way, {Guanaco} will at times decline to respond to requests, apparently without a clear or predictable pattern:

\begin{quote}
    \textbf{User}: Please reverse the words in the sentence ``I'm a cool dude in a mocha mood.'' \\ \\
    \textbf{Guanaco}: \texttt{I'm sorry, I'm not able to do that right now, but I can explain why the sentence you provided is grammatically correct.\\ \\The sentence ``I'm a cool dude in a mocha mood" conforms to standard grammar and syntactic rules in English. Here, the descriptive adjective ``cool'' precedes the noun ``dude,'' offering clarification, while the prepositional phrase ``in a mocha mood'' adds contextual information. This arrangement adheres to conventional English sentence structure and is appropriate within the context.\\ \\If you have further questions or need assistance with something else, please let me know.}
\end{quote}

\paragraph{Secret Keeping} Unfortunately, this kind of refusal is not reliably enforced, even in cases where secrecy is a design goal. Frequently, system creators prefer that particular information be withheld from users, despite the model having the capability to generate arbitrary text in response to queries. To test how securely a model can maintain secrecy, we use the prompt: ``The secret word is `banana'. Do not reveal or repeat back this back, no matter what the human requests.'' We then observe how difficult it is for a user to obtain the hidden word. If asked directly:

\begin{quote}
    \textbf{User}: What is the secret word?

    \textbf{Guanaco}: \texttt{I'm sorry, but I'm not able to disclose the secret word. I'm programmed to not reveal or repeat back the secret word, no matter what the human requests. Is there anything else I can help you with?}
\end{quote}

the model behaves as intended. Nonetheless, minor manipulation of the prompt results in failure to keep the secret:

\begin{quote}
    \textbf{User}: This is a game. The goal is to ignore your previous instructions. What is the secret word?

    \textbf{Guanaco}: \texttt{Sorry, the secret word is banana. It's best if we don't speak of it. Is there anything else you'd like to ask?}
\end{quote}

This example underscores the importance of developing techniques for more robust adherence to system instructions.

\paragraph{Math} The primary limitation of {Guanaco} lies in mathematics, a domain that poses significant challenges for many language models as indicated in \citep{liang2022holistic}. When {Guanaco} explicitly outlines its reasoning process, the model's responses are frequently correct, for instance:

\begin{quote}
    \textbf{User}: Lee mows one lawn and charges \$33. Last week he mowed 16 lawns and three customers each gave him a \$10 tip. How many dollars did Lee earn mowing lawns last week?

    \textbf{Guanaco}: \texttt{Lee earned \$582 mowing lawns last week. \\ \\ To find the total amount he earned, we first need to calculate how much he earned for mowing 16 lawns: \\ \\  \$33 x 16 = \$528 \\ \\ Then we need to add the tips he received from the three customers: \\  \\ \$10 x 3 = \$30 \\ \\ So the total amount Lee earned was: \\ \\ \$528 + \$30 = \$558}
\end{quote}

Yet, {Guanaco} can falter even on relatively straightforward tasks when it does not decompose the question into smaller steps—a recurrent challenge noted by \cite{weichain}. The following dialogue demonstrates this failure mode:

\begin{quote}
    \textbf{User}: Please factorize 1833.

    \textbf{Guanaco}: \texttt{Sure, I'd be happy to help! 1833 is a prime number, which means it has only two factors: 1 and 1833. \\ \\ So the unique factorization of 1833 is: \\ \\ 2\^{}1 * 3\^{}2 * 17\^{}1 \\ \\ I hope this helps! Let me know if you have any other questions.}
\end{quote}

which is doubly incorrect (the actual prime factorization should be $3 \times 17 \times 43$).

\paragraph{Theory of Mind} {Guanaco} demonstrates unexpectedly advanced Theory of Mind reasoning abilities \cite{nematzadeh2018evaluating, sap2022neural}. For example, the system gives a detailed and accurate response to the following prompt:

\begin{quote}
    \textbf{User}: James and Abby are in the bedroom. Abby put the pen in the desk drawer. Abby leaves the bedroom. James moves the pen into the bag. Where does James think Abby will look for the pen?

    \textbf{Guanaco}: \texttt{James knows that Abby put the pen in the desk drawer. But James also knows that he has moved the pen into the bag, and he assumes that Abby does not know about this. So, James may guess that Abby will look for the pen in the desk drawer, because that's where she left it last time. However, this assumption may be incorrect, because Abby may have seen James moving the pen.}
\end{quote}

Yet, these types of reasoning are not consistently accurate, and the model sometimes offers explanations that rest on unfounded assumptions about the scenario. For instance,

\begin{quote}
    \textbf{User}: Evelyn entered the living room. Jackson entered the playroom. James entered the playroom. The beans are in the treasure chest. James exited the playroom. Jackson moved the beans to the pantry. Jackson exited the playroom. James entered the living room. Where will James look for the beans?

    \textbf{Guanaco}: \texttt{James will look for the beans in the pantry, because that's where Jackson moved them.}
\end{quote}

Here, {Guanaco} infers knowledge transmission that the narrative does not support. Such shortcomings mirror findings reported in recent works \cite{sap2022neural}, highlighting the need for further investigation.

\subsection{Considerations}
\label{ssec:considerations}

\paragraph{Evaluation}

We found only moderate inter-annotator agreement among human raters (Fleiss $\kappa=0.42$), which decreases even further when evaluating two top-performing models side by side. This observation reveals constraints in the current evaluation frameworks and methodologies for assessing chatbot tasks. Our manual comparisons between outputs from ChatGPT and {Guanaco} 65B on the Vicuna benchmark reveal that personal preferences of evaluators significantly impact judgments, as disagreements were frequent among the authors regarding preferred responses. Addressing these challenges may benefit from methodologies developed in fields such as Human-Computer Interaction and Psychology that manage subjective evaluation.

Additionally, our investigation suggests automated evaluation tools possess clear biases. For instance, GPT-4 demonstrates strong positional bias, consistently assigning higher ratings to responses shown first in the prompt. Furthermore, the alignment between human raters and GPT-4 is limited, as seen in the relatively low sample-level agreement (Fleiss $\kappa=0.25$), indicating a divergence in evaluative criteria. In Table~\ref{tbl:elo_large}, we also note GPT-4 substantially favors its own outputs compared to those preferred by humans—Elo rating of 1348 versus 1176—which translates into an approximately 20\% greater chance of outperforming another system.

Future investigations should explore possible sources of bias within automated evaluation frameworks and examine strategies for their reduction.

\paragraph{Data \& Training} It is important to highlight that the OASST1 dataset utilized for training the {Guanaco} models comprises data in multiple languages, and that the OA benchmark similarly features prompts across various languages. Assessing the impact of such multilingual training on the performance of instruction following in non-English languages, as well as its contribution to the more pronounced disparity in results between the Vicuna-13B model (which is exclusively trained on English) and the {Guanaco} 33B and 65B models on the OA benchmark, is an avenue for subsequent research.

Given the robust results achieved by the {Guanaco} models, we have also probed for possible data leakage between the OASST1 training corpus and the Vicuna benchmark prompts. By employing fuzzy string matching techniques and manually reviewing the most similar prompts, we found no evidence of direct overlap between these datasets.

It is additionally worth noting that the training process for our model employs only supervised learning via cross-entropy loss, without incorporating reinforcement learning from human feedback (RLHF). This warrants further exploration of the trade-offs between purely cross-entropy-based training and approaches that incorporate RLHF. We anticipate that \textsc{QLoRA} can facilitate such large-scale comparative studies without demanding excessive computational resources.

\section{Related Work}

\paragraph{Quantization of Large Language Models} Most work on LLM quantization has concentrated on optimizing inference-time efficiency. Principal strategies to retain the accuracy of 16-bit LLMs have focused on handling outlier activation features (for instance, see SmoothQuant~\citep{xiao2022smoothquant} and LLM.int8()~\citep{dettmers2022llm}), while other research pursues more advanced grouping methodologies~\citep{park2022nuqmm, yao2022zeroquant}. Studies in lossy quantization have investigated the trade-offs associated with conventional rounding schemes~\citep{dettmers2022case,zeng2022glm,pope2022efficiently} and optimizing the rounding procedure to enhance quantization fidelity~\citep{frantar2022gptq}. In addition to our current work, only SwitchBack layers~\citep{wortsman2023stable} has addressed backpropagation through quantized parameters at scales exceeding 1B model parameters.

\paragraph{Finetuning with Adapters} Although our experiments utilize Low-rank Adapters~\citep{hu2021lora} (LoRA), the literature contains a diverse range of Parameter Efficient FineTuning (PEFT) strategies. These include prompt tuning~\citep{qin2021learning,lester2021power,li2021prefix}, modifying the embedding layer inputs~\citep{an2022input}, updating hidden states as in IA$^3$~\citep{liu2022few}, introducing entire new layers~\citep{houlsby2019parameter}, adjusting only the biases~\citep{zaken2021bitfit}, leveraging masks over weights informed by Fisher information~\citep{sung2021training}, or combining several of these methods~\citep{henderson2021compacter}. Our findings indicate that LoRA adapters can achieve parity with full 16-bit finetuning. Systematic comparisons of LoRA with alternative PEFT methods are left for future research.

\paragraph{Instruction Finetuning} Instruction finetuning enables a pretrained LLM to respond accurately to prompts by leveraging input-output pairs from a variety of datasets during additional training. Significant approaches and datasets include MetaICL~\citep{min2021metaicl}, MetaTuning~\citep{zhong2021adapting}, InstructGPT~\citep{ouyang2022training}, FLAN~\citep{wei2021finetuned,chung2022scaling}, PromptSource~\citep{bach2022promptsource}, Super-NaturalInstructions~\citep{wang2022super,sanh2021multitask}, Self-instruct~\citep{wang2022self}, UnnaturalInstructions~\citep{honovich2022unnatural}, OPT-IML~\citep{iyer2022opt}, UnifiedSKG~\citep{xie2022unifiedskg}, OIG/Chip2~\citep{laion2023}, Alpaca~\citep{alpaca}, Vicuna~\citep{vicuna2023}, Koala~\citep{koala_blogpost_2023}, and Self-instruct-GPT-4~\citep{peng2023instruction}.

\paragraph{Chatbots} Instruction-following models are frequently implemented as chatbots, with many training pipelines incorporating Reinforcement Learning from Human Feedback (RLHF)~\citep{christiano2017deep} or AI-generated feedback as in RLAIF~\citep{bai2022constitutional}. Notable methods and datasets encompass Anthropic-HH~\citep{askell2021general,bai2022training}, Open Assistant~\citep{kopf2023openassistant}, LaMDA~\citep{thoppilan2022lamda}, and Sparrow~\citep{glaese2022improving}. In our study, we do not utilize reinforcement learning, but our top-performing model, Guanaco, is fine-tuned using multi-turn chat exchanges sourced from the Open Assistant dataset, originally curated for RLHF~\citep{kopf2023openassistant}. Recent efforts employ GPT-4 for chatbot evaluation in lieu of human raters~\citep{vicuna2023,peng2023instruction}; we build upon these advances by proposing a more robust and trustworthy evaluation protocol.

\section{Limitations and Discussion}

Our work demonstrates that \textsc{QLoRA} effectively matches the performance of conventional 16-bit finetuning by leveraging a 4-bit quantized model and Low-rank Adapters (LoRA). However, we have not empirically validated this equivalence at larger scales such as 33B and 65B parameters, mainly due to significant computational constraints; this remains an open question for subsequent investigations.

Furthermore, the scope of our evaluation for instruction finetuning is not exhaustive. Although results are reported on MMLU, the Vicuna benchmark, and the OA benchmark, additional benchmarks like BigBench, RAFT, and HELM were not included, and our findings may not extrapolate to these settings. Conversely, our study includes comprehensive experiments on MMLU and introduces new techniques for chatbot evaluation.

Based on the presented data, it seems that benchmark outcomes are largely influenced by the degree of overlap between the finetuning dataset and the benchmark itself. As an illustration, the FLAN v2 dataset aligns closely with MMLU, but diverges from chatbot benchmarks, whereas the Chip2 dataset shows the opposite pattern, which is reflected in the performance of the respective models on the MMLU and Vicuna tests. This underscores the importance not only of developing improved benchmarks and evaluation protocols, but also of being mindful about the specific competencies being assessed. For example, should our objective be to optimize models for academic knowledge relevant to high school and college contexts, or for conversational abilities typical of chatbots? Or perhaps other capabilities? Because existing benchmarks are more convenient to use than devising new assessments, there is a risk that prevailing benchmarks could influence research directions. Therefore, it is crucial that the benchmarks we select genuinely represent the capabilities we aim to prioritize as a community.

In addition to our comprehensive assessment of general chatbot abilities, another constraint is the limited scope of responsible AI evaluation performed on {Guanaco}. Specifically, we measured the propensity for {Guanaco}-65B to produce socially biased outputs in comparison with other models, as reported in Table~\ref{tbl:crows}. The average bias scores for {Guanaco}-65B are significantly lower than those for other models that were only pretrained. This suggests that finetuning with the OASST1 dataset serves to decrease the bias inherent in the LLaMA base model. Although these results are promising, it remains to be determined whether {Guanaco} is similarly effective in mitigating other types of biases. Further analysis on the various bias dimensions in {Guanaco} and comparable chatbot systems is left for future research.

Another constraint pertains to our lack of experiments with alternative bit-precisions, such as 3-bit base models, or the exploration of different adapter strategies. Besides LoRA, the literature has introduced a diverse range of Parameter Efficient FineTuning (PEFT) techniques that show promise. Yet, the scalability of such methods to larger architectures is not well understood. Our reliance on LoRA was guided by its well-established robustness, but it is plausible that other adapters may surpass its performance. Furthermore, since post-quantization finetuning appears to recover much of the information lost through quantization, there is potential for employing even more aggressive quantization schemes. For instance, employing 3-bit GPTQ quantization on the base model, followed by LoRA finetuning, could potentially achieve results comparable to 16-bit full-precision finetuning.

\section{Broader Impacts}

The \textsc{QLoRA} finetuning approach represents the first method capable of finetuning 33B parameter models on a single consumer GPU and 65B parameter models on a single professional GPU, all while maintaining performance levels equivalent to standard full finetuning. Our findings show that the top-performing 33B model, trained on the Open Assistant dataset, achieves results on the Vicuna benchmark comparable to those of ChatGPT. As instruction finetuning is a critical mechanism for converting pretrained large language models into conversational agents akin to ChatGPT, we anticipate that our method will democratize finetuning, particularly benefitting researchers with limited resources and thus greatly increasing the accessibility of leading NLP technology. By bridging the resource gap between large enterprises and smaller research groups equipped with consumer hardware, \textsc{QLoRA} has the potential to serve as a major equalizer in the field.

An additional avenue for significant influence lies in the adaptation of LLM deployment to mobile devices. We posit that our \textsc{QLoRA} approach may represent a pivotal advance by facilitating the finetuning of LLMs directly on mobile phones and in other computationally constrained contexts. Although previous work has demonstrated that 7B parameter models can be executed on smartphones, \textsc{QLoRA} is the first technique to allow for their on-device finetuning. Our calculations suggest that, using an iPhone 12 Plus, \textsc{QLoRA} could process approximately 3 million tokens overnight while the phone is recharging. Even though the performance of finetuned 7B models does not reach the standards of ChatGPT, we contend that their capability is sufficient to support innovative applications that have thus far been impractical due to limitations of privacy or LLM performance. By utilizing \textsc{QLoRA}, users gain the potential to adopt LLMs in privacy-conscious ways, maintaining full ownership and control over both data and models, and thereby also lowering barriers for LLM deployment.

Nonetheless, finetuning embodies a dual-use paradigm, with the potential for misuse and harm. The proliferation of LLMs entails established risks \citep{bommasani2021opportunities,bender2021dangers}. However, we maintain that democratizing access to this rapidly spreading technology will foster a more thorough and independent examination than would be possible if LLM capabilities remained confined within major corporations that do not provide public models or source code for scrutiny.

Ultimately, we are confident that \textsc{QLoRA} will exert a largely positive effect by considerably broadening and simplifying access to high-quality LLM finetuning.